\section{Procedimiento de una PCR}

De manera general, la realización de una PCR requiere de los siguientes elementos:

\begin{itemize}
    \item Agua
    \item Buffer
    \item Muestra de ADN o ARN
    \begin{itemize}
        \item En caso de ser ARN se necesita una enzima retrotranscriptasa (RT)
    \end{itemize}
    \item Enzima polimerasa (Taq-polimerasa)
    \item Cofactor de la polimerasa
    \item Primers
    \item dNTP's
\end{itemize}

\subsection{Reactivos}

Primero que nada, generalmente se recomienda el uso de agua con cierta calidad de pureza, como el agua tratada con DEPC (di-etil-pirocarbonato), o el agua libre de nucleasas.

El buffer se utiliza para mantener las muestras estables.

La muestra de ADN o ADNc (complementario). Es la muestra que se desea amplificar.

La enzima polimerasa es la protagonista de la reacción. Elonga las secciones de interés de la muestra de ADN proporcionada. Puesto que los procesos de la PCR necesitan de cambios de la temperatura, llegando incluso hasta 98°C, el usar polimerasas comunes no es viable, pues sería necesario agregar la enzima cada vez que la temperatura suba, lo que aumentaría los costos del proceso, el tiempo que tomaría realizarlo y su dificultad. En 1964, Thomas D. Brock y Hudson Freezen descubrieron un organismo termófilo en las aguas termales de Yellowstone, y lo llamaron \sci{Thermus aquaticus}. En 1980, Kary Mullis, trabajando para Cetus Corporation, utilizó la polimerasa de este organismo para el nuevo procedimiento que estaba desarrollando, la PCR. Esta enzima, llamada Taq-polimerasa, podía soportar las altas temperaturas necesarias en la PCR gracias a las características del organismo de donde se aisló.

La enzima polimerasa necesita de un cofactor para poder realizar su función. El cofactor más utilizado es el \textbf{cloruro de magnesio (MgCl2)}.

Los primers son fragmentos de ADN que se unen a las cadenas de ADN, esto con el objetivo de indicarle a la polimerasa el punto de inicio para llevar a cabo la elongación.
Estos primers suelen tener un tamaño de entre 10 y 15 pares de bases.

Por último, los dNTP's son nucleotidos "sueltos" agregados al medio con los cuales la polimerasa puede trabajar y sintetizar ADN.

\subsection{Materiales}

Los materiales necesarios para realizar una PCR suelen ser pocos:

\begin{itemize}
    \item Micropipetas
    \item Tubos eppendorf (de entre 1.5ml y 0.2ml)
    \item Tubos o placas ópticas de 0.2ml
\end{itemize}

\subsection{Equipos}

Respecto a los equipos, al igual que los materiales, suelen ser pocos y siempre los mismso los que se usan para realizar una PCR

\begin{itemize}
    \item Termociclador: Es el equipo encargado de manejar los cambios de temperatura necesarios para que ocurra la PCR. Existen varios tipos de termocicladores.
        \begin{itemize}
            \item Clásico
            \item Tiempo real
            \item Digital
        \end{itemize}
    
    \item Centrífuga: Es utilizada para separar fases de la muestra.  
\end{itemize}